% ============================================================
%  INTRODUCTION GÉNÉRALE
% ============================================================
\chapter*{Introduction générale}
\addcontentsline{toc}{chapter}{Introduction générale}
\markboth{Introduction générale}{Introduction générale}

\hspace{1cm}Dans un contexte de digitalisation croissante, les \gls{pme} marocaines restent particulièrement exposées aux risques cybersécurité, sans toujours disposer des moyens de les évaluer ni de les prioriser. Le \gls{cmrpi}, à travers son expertise en formation et sensibilisation à la cybersécurité, dispose d'un référentiel de bonnes pratiques élaboré avec l'\gls{ausim}, mais celui-ci reste aujourd'hui un document statique, non opérationnalisé sous forme d'outil d'auto-diagnostic.

Le présent travail s'inscrit dans le cadre d'un stage de Projet de Fin d'Année (PFA) réalisé au sein du \gls{cmrpi} -- Espace Maroc Cyberconfiance, pour la formation d'Ingénieur d'État en \textbf{Intelligence Artificielle et Cybersécurité (IACS)} à l'École Nationale des Sciences Appliquées de Béni Mellal.

\vspace{0.5cm}
La problématique abordée concerne la conception d'une plateforme simple et accessible permettant à une \gls{pme} marocaine d'auto-évaluer sa maturité cybersécurité à partir du guide \gls{cmrpi}/\gls{ausim}, et d'obtenir des recommandations priorisées et actionnables. Pour y répondre, je propose une plateforme combinant un questionnaire structuré, un moteur de scoring automatisé fondé sur des référentiels normatifs reconnus, et la génération de rapports personnalisés.

\vspace{0.5cm}
Ce rapport est organisé en trois chapitres principaux :

\begin{itemize}
  \item \textbf{Chapitre I -- Contexte général et problématique :} présentation du \gls{cmrpi}, de la problématique et du cahier des charges du stage.
  \item \textbf{Chapitre II -- État de l'art :} analyse des référentiels normatifs mobilisés (guide \gls{cmrpi}/\gls{ausim}, \gls{iso27001}, \gls{nistcsf}, loi 09-08) et justification académique de la méthodologie retenue.
  \item \textbf{Chapitre III -- Conception et réalisation :} description de l'architecture proposée, des choix technologiques, de la méthodologie de scoring, de l'implémentation et des résultats obtenus.
\end{itemize}
